\chapter{Conclusioni}
\label{chap:final-conclusions}

Il presente progetto di semestre, svolto in collaborazione con il
Politecnico di Milano, ha affrontato il problema della
\textbf{predizione di parametri di celle Litio-Ione} a partire da un
dataset di spettroscopia di impedenza galvanostatica composto da
\num{9805} misurazioni e 200 curve di Nyquist complete.

\section*{Risultati ottenuti}

I tre task definiti nel \cref{chap:introduction} sono stati tutti
risolti con metriche superiori alla soglia operativa di interesse
($R^{2} \geq 0.95$ per la regressione, accuracy $\geq 0.95$ per la
classificazione):
\begin{itemize}
  \item \textbf{Task 1 -- Leave-One-Out:} ricostruzione di una coppia
        $(Aging, Temperatura)$ esclusa dal training con
        $R^{2} = 0.991$ tramite \emph{KNN distance-weighted}, MSE di
        $0.0079\,\milliohm^2$ e tempo di addestramento sotto il
        secondo;
  \item \textbf{Task 2 -- Young/Old Classification:} classificazione
        binaria delle 8 curve di Nyquist (una per temperatura) di una
        coppia $(Aging, SOC)$ esclusa dal training, con accuracy
        $1.000$ sul hold-out di default $(Aging = 4, SOC = 3)$ tramite
        \emph{SVM-RBF}; il modello satura agli estremi del dominio
        (Aging $\in \{0, 4\}$) e si comporta come un vero stress test
        alle coppie \emph{borderline} attorno alla soglia young/old
        (Aging $\in \{2, 3\}$), dove l'accuracy oscilla fra $0.50$
        e $1.00$ in funzione del SOC scelto;
  \item \textbf{Task 3 -- Aging Interpolation:} predizione di un
        intero livello di Aging mai osservato con $R^{2} = 0.986$
        tramite \emph{KNN distance-weighted}, con degrado controllato
        ai bordi della finestra termica.
\end{itemize}

\section*{Impatto operativo stimato}

Sul piano operativo, i numeri suggeriscono che è
ragionevole \textbf{evitare la misura} di una coppia (Aging,
Temperatura) ogni 40, ovvero un risparmio del $\sim$2.5\,\% sul tempo
di laboratorio per ogni cella caratterizzata, mantenendo l'errore di
predizione su $\mathrm{Re}(Z)$ e $\mathrm{Im}(Z)$ ben dentro la
dispersione fisiologica delle misure ($\sim$\SI{1}{\milli\ohm}). Per il
Task 3 il risparmio potenziale è ancora più rilevante: predire un
intero stadio di invecchiamento permette di saltare un'intera campagna
di degrado accelerato, dell'ordine di decine di ore di laboratorio.

\section*{Generalizzazione e contributo metodologico}

L'esperienza ha confermato tre principi che riteniamo riusabili in
contesti simili (qualsiasi problema di regressione/classificazione su
dati misurati a costo non nullo):
\begin{enumerate}[leftmargin=*]
  \item lo \textbf{splitter di cross-validation} deve imitare la
        procedura di valutazione fuori-sample del task -- usare una
        CV row-wise quando il test è un gruppo intero produce numeri
        sistematicamente ottimistici;
  \item il \textbf{feature engineering motivato fisicamente}
        (Arrhenius, log-frequenza, interazioni $X \times T$) batte
        ensemble ciechi più sofisticati;
  \item l'\textbf{infrastruttura MLOps} (configurazione centralizzata,
        tracking, test, drift detection) ha un costo fisso modesto e
        si ripaga la prima volta che bisogna riprodurre un risultato a
        distanza di settimane.
\end{enumerate}

\section*{Prospettive future}

Le direzioni più promettenti, in ordine di priorità, sono:
\begin{enumerate}[leftmargin=*]
  \item aggiungere \textbf{intervalli di confidenza bootstrap} alle
        metriche di benchmark, per rendere statisticamente difendibili
        i confronti fra modelli (~30 righe di codice in
        \codefile{src/evaluation/metrics.py});
  \item chiudere il ciclo di \textbf{continual learning} legando i
        report di drift a un retraining automatico orchestrato;
  \item esplorare \textbf{modelli ibridi con prior fisico} (Arrhenius
        parametrizzato, polinomi di basso grado in Aging) per
        affrontare l'estrapolazione oltre il dominio osservato;
  \item versionare il dataset con \textbf{DVC} e migrare l'experiment
        tracking su un backend MLflow remoto, per supportare
        collaborazione fra più sviluppatori.
\end{enumerate}

In sintesi, il progetto ha consegnato i tre notebook richiesti e una
piattaforma MLOps funzionante che li espone, traccia e monitora. I
limiti riconosciuti nel \cref{chap:conclusions} sono noti e
quantificati, e tracciano una roadmap operativa per il lavoro futuro.
